% !TeX root = ../main.tex
\chapter{Supplementary Specifications}
\label{appendix:supplementary-specifications}

This appendix contains compact reference material that supports the
analysis and design chapters. Detailed diagrams are separated into
dedicated appendices by diagram type so that Chapter~\ref{chap:phan-tich-thiet-ke}
can focus on the main architectural narrative.

\section{Appendix Organization}
\label{appendix:organization}

\begin{table}[H]
    \centering
    \caption{Appendix organization.}
    \label{tab:appendix-organization}
    \begin{tabular}{|p{2.5cm}|p{4.2cm}|p{6.7cm}|}
        \hline
        \textbf{Appendix} & \textbf{Content} & \textbf{Purpose} \\
        \hline
        Appendix~\ref{appendix:supplementary-specifications}
            & Supplementary specifications
            & Detailed specs for secondary use cases plus database schema summary. \\
        \hline
        Appendix~\ref{appendix:sequence-diagrams}
            & Detailed sequence diagrams
            & Step-by-step interaction flows that are too detailed for Chapter 4. \\
        \hline
        Appendix~\ref{appendix:activity-diagrams}
            & Detailed activity diagrams
            & Decision points and parallel branches of major workflows. \\
        \hline
        Appendix~\ref{appendix:class-diagrams}
            & Detailed class diagrams
            & Additional bounded-context diagrams supporting the domain model. \\
        \hline
        Appendix~\ref{appendix:state-machines}
            & Detailed state machines
            & Operational and review lifecycle states separated from the main design flow. \\
        \hline
        Appendix~\ref{appendix:security-details}
            & Detailed security controls
            & Full OWASP LLM risk analysis and mitigation notes. \\
        \hline
        Appendix~\ref{appendix:ui-details}
            & Detailed interface design
            & Sitemap and screen-level UI descriptions. \\
        \hline
        Appendix~\ref{appendix:operational-details}
            & Operational and integration details
            & Proto contracts, observability fields, and deployment topology details. \\
        \hline
    \end{tabular}
\end{table}

\section{Detailed Specifications of Secondary Use Cases}
\label{appendix:use-case-reference}

The use case diagram and detailed specifications of the main use cases
are in
Section~\ref{sec:use-case-diagram} of
Chapter~\ref{chap:khao-sat-phan-tich-he-thong} (specifically
Section~\ref{subsec:main-use-case-specs}). This appendix collects the
detailed specifications of the remaining eleven secondary use cases
using the same template (Primary actor, Preconditions, Main flow,
Postconditions, Alternate flows). They are placed here because their
flows are operational or auxiliary rather than load-bearing for the
main thesis defense.

\subsection{UC-IN-02 -- Monitor Processing}

\begin{longtable}{|p{3.2cm}|p{10.3cm}|}
\hline
\textbf{Field} & \textbf{Specification} \\
\hline
\endhead
Primary actor & Instructor (AC-01) \\
\hline
Preconditions & At least one document or video is in a non-terminal
    pipeline state (\texttt{QUEUED} through \texttt{EMBEDDING}). \\
\hline
Main flow &
    1.\ Instructor opens the upload list page in \texttt{web-bff}.
    2.\ \texttt{web-bff} subscribes to Redis Pub/Sub and forwards
    progress events to the browser via SSE.
    3.\ Each row updates in place as pipeline state changes
    (\texttt{PARSING} $\to$ \texttt{CHUNKING} $\to$ \texttt{EMBEDDING}
    $\to$ \texttt{INDEXED}). \\
\hline
Postconditions & UI shows the latest pipeline state; instructor can
    retry items in \texttt{ERROR} state. \\
\hline
Alternate flows &
    \emph{A1} - SSE connection drops: browser auto-reconnects; state
    backfilled from the next snapshot.
    \emph{A2} - Item in \texttt{ERROR}: instructor clicks Retry; a
    new job is enqueued. \\
\hline
\end{longtable}

\subsection{UC-IN-05 -- View Class Analytics}

\begin{longtable}{|p{3.2cm}|p{10.3cm}|}
\hline
\textbf{Field} & \textbf{Specification} \\
\hline
\endhead
Primary actor & Instructor (AC-01) \\
\hline
Preconditions & The course has at least one quiz attempt and one
    learning event from a learner. \\
\hline
Main flow &
    1.\ Instructor opens \texttt{/manage/analytics}.
    2.\ \texttt{core-api} aggregates \texttt{quiz\_attempts} and
    \texttt{learning\_events} grouped by Learning Outcome.
    3.\ Dashboard shows: active learners, completion rate, average
    quiz score, top 5 weakest LOs. \\
\hline
Postconditions & No state change. \\
\hline
Alternate flows &
    \emph{A1} - No data yet: dashboard shows ``insufficient data''
    placeholder with a guide to onboard learners. \\
\hline
\end{longtable}

\subsection{UC-IN-06 -- Manage Learning Resources}

\begin{longtable}{|p{3.2cm}|p{10.3cm}|}
\hline
\textbf{Field} & \textbf{Specification} \\
\hline
\endhead
Primary actor & Instructor (AC-01) \\
\hline
Preconditions & Instructor launched the tool from Canvas with role
    \texttt{Instructor}; at least one document or video belongs to the
    selected course. \\
\hline
Main flow &
    1.\ Instructor opens the resource library in \texttt{web-bff}.
    2.\ \texttt{core-api} returns documents/videos with processing
    status, publish status, owner, and derived artifact counts.
    3.\ Instructor selects a resource and chooses Rename, Rerun
    pipeline, or Delete.
    4.\ For Rename, \texttt{core-api} updates metadata and writes an
    audit record.
    5.\ For Rerun, \texttt{core-api} enqueues a new processing job
    while preserving the old result until the new result is ready.
    6.\ For Delete, \texttt{core-api} soft-deletes the resource and
    publishes \texttt{DocumentDeleted}/\texttt{VideoDeleted} for
    cleanup across storage, vector index, and graph stores. \\
\hline
Postconditions & Resource metadata is updated, the pipeline is rerun,
    or derived artifacts are marked for idempotent cleanup. Audit logs
    preserve the instructor action. \\
\hline
Alternate flows &
    \emph{A1} - Resource has published content: the UI requires a
    second confirmation before deletion.
    \emph{A2} - Delete is repeated: \texttt{core-api} returns success
    without creating duplicate cleanup effects. \\
\hline
\end{longtable}

\subsection{UC-LE-02 -- Review with Flashcards}

\begin{longtable}{|p{3.2cm}|p{10.3cm}|}
\hline
\textbf{Field} & \textbf{Specification} \\
\hline
\endhead
Primary actor & Learner (AC-02) \\
\hline
Preconditions & At least one \texttt{PUBLISHED} card in the lesson. \\
\hline
Main flow &
    1.\ Learner opens the lesson and toggles flashcard mode.
    2.\ Cards are presented as flippable front/back (title /
    key\_insight).
    3.\ After flipping, learner selects ``Mastered'' or ``Review
    again''; \texttt{flashcard\_reviews} is upserted with new
    \texttt{state}, \texttt{last\_seen\_at},
    \texttt{next\_due\_at}. \\
\hline
Postconditions & Per-learner spaced-repetition state is updated. \\
\hline
Alternate flows &
    \emph{A1} - No published cards: UI shows the standard lesson view
    instead of flashcard mode. \\
\hline
\end{longtable}

\subsection{UC-LE-05 -- View Personal Progress}

\begin{longtable}{|p{3.2cm}|p{10.3cm}|}
\hline
\textbf{Field} & \textbf{Specification} \\
\hline
\endhead
Primary actor & Learner (AC-02) \\
\hline
Preconditions & Learner has interacted with at least one published
    lesson. \\
\hline
Main flow &
    1.\ Learner opens \texttt{/learn/progress}.
    2.\ \texttt{core-api} aggregates the learner's
    \texttt{learning\_events} and \texttt{quiz\_attempts} grouped by
    LO.
    3.\ Dashboard shows: percentage of cards viewed, quiz score
    trend, top 3 weakest LOs to review. \\
\hline
Postconditions & No state change. \\
\hline
Alternate flows & None. \\
\hline
\end{longtable}

\subsection{UC-LE-06 -- Search Course Content}

\begin{longtable}{|p{3.2cm}|p{10.3cm}|}
\hline
\textbf{Field} & \textbf{Specification} \\
\hline
\endhead
Primary actor & Learner (AC-02) \\
\hline
Preconditions & Learner launched the tool from Canvas with role
    \texttt{Learner}; the course has at least one published card,
    chunk, or transcript segment indexed for retrieval. \\
\hline
Main flow &
    1.\ Learner types a keyword or natural-language query in the
    course search box.
    2.\ \texttt{web-bff} forwards the query with course and role
    context to \texttt{chat-service} or \texttt{core-api}.
    3.\ The retrieval layer applies metadata filters
    (\texttt{course\_id}, publish status, role) and runs hybrid search
    over Qdrant.
    4.\ Results are returned as cards/chunks/transcript segments with
    source document, page, heading, or video timestamp. \\
\hline
Postconditions & No learning state is changed unless the learner opens
    a result; search logs can be stored for retrieval evaluation. \\
\hline
Alternate flows &
    \emph{A1} - No result: UI shows an empty state and suggests asking
    the AI Tutor.
    \emph{A2} - Query is outside the course scope: only published
    course-scoped results are returned; no external knowledge is used. \\
\hline
\end{longtable}

\subsection{UC-LMS-02 -- Deep Linking}

\begin{longtable}{|p{3.2cm}|p{10.3cm}|}
\hline
\textbf{Field} & \textbf{Specification} \\
\hline
\endhead
Primary actor & Instructor (AC-01) \\
\hline
Secondary actor & Canvas LMS (AC-03) \\
\hline
Preconditions & At least one lesson with \texttt{PUBLISHED} content;
    Canvas tool registration has the Deep Linking message type
    enabled. \\
\hline
Main flow &
    1.\ Instructor adds an external tool in Canvas; Canvas POSTs a
    Deep Linking request JWT to \texttt{/lti/deep-link-select}.
    2.\ \texttt{web-bff} validates the JWT, lets the instructor pick a
    target (lesson, card, quiz\_set, chat, or video).
    3.\ \texttt{web-bff} constructs an
    \texttt{LtiResourceLinkRequest} JWT pointing at the chosen target
    and POSTs it back to Canvas \texttt{deep\_link\_return\_url}.
    4.\ \texttt{core-api} persists an \texttt{lti\_resource\_links}
    row mapping \texttt{lms\_resource\_link\_id} $\to$
    \texttt{target\_kind} + \texttt{target\_id}. \\
\hline
Postconditions & Canvas now has an activity that, when launched,
    routes the learner straight to the chosen target. \\
\hline
Alternate flows &
    \emph{A1} - Instructor cancels: \texttt{web-bff} returns to
    Canvas with an empty content-items array; no row is written. \\
\hline
\end{longtable}

\subsection{UC-LMS-03 -- Score Synchronization (AGS)}

\begin{longtable}{|p{3.2cm}|p{10.3cm}|}
\hline
\textbf{Field} & \textbf{Specification} \\
\hline
\endhead
Primary actor & System (triggered by UC-LE-03 quiz submission) \\
\hline
Secondary actor & Canvas LMS (AC-03) \\
\hline
Preconditions & The quiz attempt exists; an
    \texttt{lti\_resource\_links} row references the attempt's lesson;
    a corresponding \texttt{lti\_line\_items} row exists. \\
\hline
Main flow &
    1.\ \texttt{core-api} fetches an OAuth 2.0 access token from
    Canvas (client\_credentials with private-key JWT assertion).
    2.\ POST to \texttt{lti\_line\_items.lms\_line\_item\_url}
    \texttt{/scores} with \texttt{userId}, \texttt{scoreGiven},
    \texttt{scoreMaximum}, \texttt{activityProgress},
    \texttt{gradingProgress}.
    3.\ On 2xx: update \texttt{quiz\_attempts.ags\_status =
    POSTED} and \texttt{ags\_posted\_at}. \\
\hline
Postconditions & Canvas Gradebook reflects the new score within
    $\leq 5$ seconds. \\
\hline
Alternate flows &
    \emph{A1} - Non-2xx response: \texttt{ags\_status = FAILED} with
    error captured; a retry job runs with exponential backoff.
    \emph{A2} - Token cache miss/expired: refresh OAuth token then
    retry. \\
\hline
\end{longtable}

\subsection{UC-AD-01 -- Manage AI Limits}

\begin{longtable}{|p{3.2cm}|p{10.3cm}|}
\hline
\textbf{Field} & \textbf{Specification} \\
\hline
\endhead
Primary actor & System Administrator (AC-05) \\
\hline
Preconditions & Administrator authenticated via admin-scope JWT. \\
\hline
Main flow &
    1.\ Administrator opens \texttt{/admin/quota}.
    2.\ \texttt{core-api} lists current entries in
    \texttt{user\_llm\_quota} and \texttt{scope\_llm\_quota}.
    3.\ Administrator edits limits (daily / monthly tokens, per-user
    or per-scope) and saves.
    4.\ \texttt{core-api} validates and persists; rate-limit layer
    picks up new limits on the next check. \\
\hline
Postconditions & Updated limits take effect; previously blocked users
    can be unblocked manually if needed. \\
\hline
Alternate flows &
    \emph{A1} - Limit raised above hard cap: rejected with HTTP 422.
    \emph{A2} - Bulk import via CSV: validated in a single
    transaction. \\
\hline
\end{longtable}

\subsection{UC-AD-02 -- Monitor System}

\begin{longtable}{|p{3.2cm}|p{10.3cm}|}
\hline
\textbf{Field} & \textbf{Specification} \\
\hline
\endhead
Primary actor & System Administrator (AC-05) \\
\hline
Preconditions & Observability stack (Prometheus / Grafana / Loki /
    Tempo) running and scraping all runtimes. \\
\hline
Main flow &
    1.\ Administrator opens the Grafana subdomain.
    2.\ Dashboards display: active jobs, error rate (24h), latency
    p95, queue depth, accumulated LLM cost.
    3.\ Administrator can open Bull-board to inspect failed jobs and
    requeue. \\
\hline
Postconditions & No state change in business data; failed jobs may
    be requeued. \\
\hline
Alternate flows &
    \emph{A1} - Threshold-based alert fired: alert delivered to
    Discord/Email; administrator follows the on-call runbook. \\
\hline
\end{longtable}

\subsection{UC-AD-03 -- Configure LMS Integration}

\begin{longtable}{|p{3.2cm}|p{10.3cm}|}
\hline
\textbf{Field} & \textbf{Specification} \\
\hline
\endhead
Primary actor & System Administrator (AC-05) \\
\hline
Secondary actor & Canvas LMS (AC-03) \\
\hline
Preconditions & Administrator has access to the AI admin interface and
    Canvas Developer Key settings. \\
\hline
Main flow &
    1.\ Administrator opens \texttt{/admin/lti-settings}.
    2.\ The system shows the required Canvas configuration values:
    issuer, client ID, deployment ID, JWKS URL, login URL, launch URL,
    redirect URLs, Deep Linking URL, and AGS scopes.
    3.\ Administrator enters or updates the values obtained from
    Canvas Developer Key configuration.
    4.\ \texttt{core-api} validates required fields, stores encrypted
    secrets, and tests JWKS retrieval.
    5.\ Administrator runs a test launch from Canvas; the system
    verifies issuer, audience, nonce, deployment ID, and role claims. \\
\hline
Postconditions & Canvas and the AI system are mutually configured for
    LTI launch, Deep Linking, and AGS passback. \\
\hline
Alternate flows &
    \emph{A1} - JWKS or deployment validation fails: configuration is
    saved as inactive and the UI shows the failed check.
    \emph{A2} - AGS scope missing: launch can still work, but grade
    synchronization is disabled until the scope is granted. \\
\hline
\end{longtable}

\section{Database Schema Reference}
\label{appendix:db-schema}

The complete PostgreSQL initialization DDL is separated into
\texttt{report/lvtn\_en/appendix/db\_schema.sql} to keep the report
readable while preserving an implementation-ready schema artifact. The
file includes:

\begin{itemize}
    \item \textbf{31 tables for Phase 1:} identity and LMS integration
          (lms\_user\_mappings, courses, lms\_course\_ref,
          lti\_resource\_links, lti\_line\_items, course\_memberships),
          curriculum (chapters, learning\_outcomes with academic-year
          versioning, assessments, lo\_assessments), document
          ingestion (documents, chunks, chunk\_lo\_mappings), video
          pipeline (videos, video\_segments, transcript\_segments),
          generated content (lessons, lesson\_cards,
          card\_video\_attachments, quiz\_items, flashcard\_reviews),
          analytics and audit (quiz\_attempts, chat\_messages,
          learning\_events, review\_audit\_logs, llm\_usage\_logs,
          user\_llm\_quota, scope\_llm\_quota), asynchronous outbox,
          content generation requests, and engine cache
          (content\_generation\_requests, outbox\_events,
          youtube\_search\_cache).
    \item \textbf{Three Phase-2 extension tables for adaptive learning:}
          \texttt{lo\_mastery\_states} stores the four BKT parameters
          ($p_{\text{known}}, p_{\text{learn}}, p_{\text{guess}},
          p_{\text{slip}}$) per user and LO;
          \texttt{recommendation\_logs} records bandit-generated
          recommendations (Thompson sampling or LinUCB) together with the
          posterior distribution; \texttt{learning\_paths} persists the LO
          sequence selected by beam search as the personalized path.
    \item \textbf{CHECK constraints for state enums:} the content-review
          states follow the lifecycle in
          Appendix~\ref{appendix:state-content-review}; the document/video
          processing states follow Section~\ref{sec:state-machine-design}.
    \item \textbf{Partial indexes:} \texttt{idx\_outbox\_pending\_poller}
          \texttt{WHERE status = 'PENDING'} supports the outbox polling
          daemon; \texttt{idx\_lesson\_cards\_review}
          \texttt{WHERE status IN ('GENERATED\_DRAFT','REVIEWING',
          'CHANGES\_REQUESTED')} supports the instructor review inbox.
    \item \textbf{Views:} \texttt{v\_current\_learning\_outcomes} filters the
          current LO version; \texttt{v\_lo\_mastery\_summary} aggregates
          mastery data for the instructor dashboard in Phase 2.
\end{itemize}

\subsection*{Schema Design Notes}

\textbf{Versioning Learning Outcomes.} The table
\texttt{learning\_outcomes} includes \texttt{academic\_year} and
\texttt{version} to track changes to LOs by academic year. This novelty
requirement allows the Knowledge Graph to represent curriculum evolution.
The \texttt{is\_current} flag marks the version used in RAG retrieval,
while old versions are kept for tracing learning history by cohort.

\textbf{Denormalized \texttt{course\_id}.} Tables frequently filtered by
course -- \texttt{chunks}, \texttt{video\_segments},
\texttt{lesson\_cards}, \texttt{quiz\_items}, \texttt{quiz\_attempts},
\texttt{chat\_messages} -- carry \texttt{course\_id} even when it could
be derived by traversal. The purpose is to reduce RAG queries to one
table and synchronize directly into Qdrant payload for hard tenant
filtering.

\textbf{Status enum consistent with State Machine.} The
\texttt{lesson\_cards} and \texttt{quiz\_items} tables use CHECK
constraints with seven states exactly matching the state machine in
Section~\ref{sec:state-machine-design}
(\texttt{GENERATED\_DRAFT} $\to$ \texttt{REVIEWING} $\to$
\texttt{CHANGES\_REQUESTED} $\to$ \texttt{APPROVED} $\to$
\texttt{PUBLISHED} $\to$ \texttt{UNPUBLISHED} $\to$ \texttt{ARCHIVED}).
Partial indexes on status allow review-inbox and learner-serving queries
to avoid full table scans.

\textbf{Time-ordered primary keys (UUIDv7).} High-write tables --
\texttt{chunks}, \texttt{learning\_events}, \texttt{chat\_messages},
\texttt{outbox\_events} -- use UUIDv7 as primary keys. PostgreSQL 18+
uses the built-in \texttt{uuidv7()} function; older versions use the
\texttt{pg\_uuidv7} extension. The embedded timestamp makes B-Tree index
inserts near-sequential, avoiding index-page fragmentation and improving
bulk insert speed.

\textbf{Outbox pattern.} Every change that must be synchronized to Neo4j
or downstream services is written to \texttt{outbox\_events} in the same
transaction as the business action. The table includes
\texttt{aggregate\_type}, \texttt{aggregate\_id},
\texttt{processed\_at}, and \texttt{last\_error} for synchronization
tracking and debugging failed jobs. The Outbox Poller daemon in
\texttt{core-api} periodically reads this table and pushes jobs to
BullMQ; the appropriate worker (\texttt{document-worker} for indexing and
graph sync, \texttt{enrichment-worker} for LLM tasks) consumes the job and
projects it to Neo4j or emits downstream events, ensuring eventual
consistency.

\textbf{Auth delegation.} The system does not store user passwords. User
authentication is delegated to Canvas through LTI 1.3. An internal user is
only \texttt{(internal\_user\_id, display\_name, role)} mapped from
\texttt{lms\_user\_mappings}.

The Qdrant collection design (dense and sparse named vectors, payload
indexes, hard filters by \texttt{course\_id}, \texttt{source\_type},
\texttt{lo\_ids}, \texttt{language}, and \texttt{visibility}) is
summarized in Section~\ref{subsec:qdrant-vector-schema}. The Neo4j
constraints and indexes are summarized in
Section~\ref{subsec:neo4j-knowledge-graph}.
